\documentclass[english]{article}
\usepackage{amsmath}
\usepackage{ae,aecompl}
\usepackage[T1]{fontenc}
\usepackage[latin9]{inputenc}
\usepackage{textcomp}
\usepackage{graphicx}
\usepackage{eurosym}
\usepackage[hidelinks]{hyperref}
\usepackage{float}
\usepackage{fancyhdr}
\usepackage{tabto}
\usepackage{verbatim}
\usepackage{etoolbox}
\usepackage[table]{xcolor}
\usepackage{booktabs}
\usepackage{wasysym}
\usepackage{multirow}
\usepackage{fancyhdr}
\usepackage{pdflscape}
\usepackage[scaled]{beramono}
\usepackage{listings}
\usepackage{csquotes}
\usepackage{comment}
\lstset{
	breaklines=true,
	}
\pagestyle{fancy}
\lhead{\powerenjoy I.T.P.D.}
\makeatletter
\makeatother
\usepackage{babel}

\newcommand\fnurl[2]{%
	\href{#2}{#1}\footnote{\url{#2}}%
}
\newcommand{\rent}{\textit{rent }}
\newcommand{\carsharing}{\textit {car sharing }}
\newcommand{\powerenjoy}{\textit{PowerEnjoy }}
\newcommand{\registereduser}{\textit {registered user }}
\newcommand{\registeredusers}{\textit {registered users }}
\newcommand{\powerenjoyuser}{\textit{PowerEnjoy User }}
\newcommand{\staff}{\textit{staff }}
\newcommand{\service}{\textit{service }}
\newcommand{\services }{\textit{services }}
\newcommand{\safearea}{\textit{safe area }}
\newcommand{\safeareas}{\textit{safe areas }}
\newcommand{\powerplug}{\textit{power plug }}
\newcommand{\powerplugs}{\textit{power plugs }}
\newcommand{\reservation}{\textit{reservation }}
\newcommand{\stand}{\textit{stand }}
\newcommand{\fieldstaff}{\textit{field staff }}
\newcommand{\resevation}{\textit{reservation }}
\newcommand{\stopover}{\textit{stopover }}
\newcommand{\personalpin}{\textit{personal pin }}
\newcommand{\guest}{\textit{guest }}
\newcommand{\thirdparty}{\textit{third party developer }}
\newcommand{\powerplugslot}{\textit{power plug slot }}
\newcommand{\moneysaving}{\textit{money saving }}
\begin{document}
\lstset{language=Java} 

	\begin{figure}
		\centering
		\includegraphics[scale=0.5]{logo.pdf} 
	\end{figure}


	\title{PowerEnJoy\\
	Integration Test Plan Document\\
	}

	\date{A.A 2016/2017}
	
	\author{Erba Alessandro\\
	 Leveni Filippo\\
	 Lodi Luca}
	
	\maketitle
	\pagebreak{}
	\tableofcontents{} \pagebreak{}
	
\section{IterateSectionWidget}
	\subsection{Functional role}
	\subsection{Application of the checklist}
\begin{itemize}
	\item[23.]{The javadoc is not used as shows at this link:\\
https://ci.apache.org/projects/ofbiz/site/javadocs/org/apache/ofbiz/widget/model/IterateSectionWidget.html}
	\item[25.]{line 55
		\begin{lstlisting}
@SuppressWarnings("serial")
		\end{lstlisting}
Between class documentation comment and class statement must not be any code.
}
\end{itemize}
\subsubsection*{.IterateSectionWidget}
\begin{itemize}
		\item[33.]{ 
			\begin{itemize}
				\item{line 77
						\begin{lstlisting}
String entryName = iterateSectionElement.getAttribute("entry");
						\end{lstlisting} must be moved to the top}
				\item{line 82
						\begin{lstlisting}
String keyName = iterateSectionElement.getAttribute("key");
						\end{lstlisting} must be moved to the top}
				\item{line 89
						\begin{lstlisting}
int viewSize = DEFAULT_PAGE_SIZE;
						\end{lstlisting} must be moved to the top}
				\item{line 90
						\begin{lstlisting}
String viewSizeStr = iterateSectionElement.getAttribute("view-size");
						\end{lstlisting} must be moved to the top}
				\item{line 95
						\begin{lstlisting}
List<? extends Element> childElementList = UtilXml.childElementList(iterateSectionElement);
						\end{lstlisting} must be moved to the top}
			\end{itemize}
}
		\item[39.]{ 
			\begin{itemize}
				\item{line 95
						\begin{lstlisting}
List<? extends Element> childElementList = UtilXml.childElementList(iterateSectionElement);
						\end{lstlisting} must be declared using a constructor.}
			\end{itemize}
}
\end{itemize}
\subsubsection*{.getSectionList}	
\subsubsection*{.renderWidgetString}
\begin{itemize}
	\item[27.]{The method is too long (90 lines), it should be reduced.
}
		\item[33.]{ 
			\begin{itemize}
				\item{line 132
						\begin{lstlisting}
List<?> theList = null;
						\end{lstlisting} must be moved to the top}
				\item{line 146
						\begin{lstlisting}
int startPageNumber = WidgetWorker.getPaginatorNumber(context);
						\end{lstlisting} must be moved to the top}
				\item{line 171
						\begin{lstlisting}
Iterator<?> iter = theList.iterator();
						\end{lstlisting} must be moved to the top}
				\item{line 172
						\begin{lstlisting}
int itemIndex = -1;
						\end{lstlisting} must be moved to the top}
				\item{line 173
						\begin{lstlisting}
int iterateIndex = 0;
						\end{lstlisting} must be moved to the top}
				\item{line 179
						\begin{lstlisting}
Object item = iter.next();
						\end{lstlisting} must be moved to the top of the "while" block}
			\end{itemize}
		}
		\item[39.]{ 
			\begin{itemize}
				\item{line 132
						\begin{lstlisting}
List<?> theList = null;
						\end{lstlisting} must be declared using a constructor.}
			\end{itemize}
}
	\item[40.]{ line 128
		\begin{lstlisting}
if (obj == null)
		\end{lstlisting} must be replaced using
		\begin{lstlisting}
if (obj.equals(null))
		\end{lstlisting}}
	\item[42.]{line 129
		\begin{lstlisting}
Debug.logError(``No object found for listName:''+ listNameExdr.toString(), module);
		\end{lstlisting}
		the description of the log could be more precise because in order to make more comprehensible the problem.}
	\item[42.]{line 142.		\begin{lstlisting} 
Debug.logError(``Object not list or map type'', module);
		\end{lstlisting}
		the description of the log could be more precise because in order to make more comprehensible the problem.}
	\item[42.]{line 218
		\begin{lstlisting} 
 Debug.logError(e, module);
		\end{lstlisting}
		the description of the log could be more precise because in order to make more comprehensible the problem.}
	
\end{itemize}

\subsubsection*{.getPaginateTarget}
\subsubsection*{.getPaginate}
\subsubsection*{.getViewSize}
\subsubsection*{.renderNextPrev}
	\begin{itemize}
	\item[27.]{The method is too long (100 lines), it should be reduced.
}	
		\item[33.]{ 
			\begin{itemize}
				\item{line 249
						\begin{lstlisting}
Map<String, Object> inputFields = UtilGenerics.checkMap(context.get("requestParameters"));
						\end{lstlisting} must be moved to the top}
				\item{line 250
						\begin{lstlisting}
Map<String, Object> queryStringMap = UtilGenerics.toMap(context.get("queryStringMap"));
						\end{lstlisting} must be moved to the top}
				\item{line 255
						\begin{lstlisting}
String queryString = UtilHttp.urlEncodeArgs(inputFields);
						\end{lstlisting} must be moved to the top}
				\item{line 256
						\begin{lstlisting}
int paginatorNumber = WidgetWorker.getPaginatorNumber(context);
						\end{lstlisting} must be moved to the top}
				\item{line 267
						\begin{lstlisting}
int viewIndex = -1;
						\end{lstlisting} must be moved to the top}
				\item{line 274
						\begin{lstlisting}
int viewSize = -1;
						\end{lstlisting} must be moved to the top}
				\item{line 298
						\begin{lstlisting}
int lowIndex = viewIndex * viewSize;
						\end{lstlisting} must be moved to the top}
				\item{line 299
						\begin{lstlisting}
int highIndex = (viewIndex + 1) * viewSize;
						\end{lstlisting} must be moved to the top}
				\item{line 305
						\begin{lstlisting}
HttpServletRequest request = (HttpServletRequest) context.get("request");
						\end{lstlisting} must be moved to the top}
				\item{line 306
						\begin{lstlisting}
HttpServletResponse response = (HttpServletResponse) context.get("response");
						\end{lstlisting} must be moved to the top}
				\item{line 308
						\begin{lstlisting}
ServletContext ctx = (ServletContext) request.getAttribute("servletContext");
						\end{lstlisting} must be moved to the top}
				\item{line 309
						\begin{lstlisting}
RequestHandler rh = (RequestHandler) ctx.getAttribute("_REQUEST_HANDLER_");
						\end{lstlisting} must be moved to the top}
				\item{line 317
						\begin{lstlisting}
StringBuilder linkText = new StringBuilder(targetService);
						\end{lstlisting} must be moved to the top of the "if" block}
				\item{line 328
						\begin{lstlisting}
String previous = UtilProperties.getMessage("CommonUiLabels", "CommonPrevious", (Locale) context.get("locale"));
						\end{lstlisting} must be moved to the top of the "if" block}
				\item{line 339
						\begin{lstlisting}
StringBuilder linkText = new StringBuilder(targetService);
						\end{lstlisting} must be moved to the top of the "if" block}
				\item{line 349
						\begin{lstlisting}
String next = UtilProperties.getMessage("CommonUiLabels", "CommonNext", (Locale) context.get("locale"));
						\end{lstlisting} must be moved to the top of the "if" block}
			\end{itemize}
		}
	\item[40.]{ line 245 
		\begin{lstlisting}
if (targetService = = null) 
		\end{lstlisting} must be replaced using
		\begin{lstlisting}
if (targetService.equals(null))
		\end{lstlisting}}
	
	
	 \item[42.]{line 263
		\begin{lstlisting} 
 Debug.logWarning(``TargetService is empty.'', module);
		\end{lstlisting}
		the description of the log could be more precise because in order to make more comprehensible the problem.}
\end{itemize}
\subsubsection*{.accept}
\subsubsection*{.getListNameExdr}
\subsubsection*{.getEntryNameExdr}
\subsubsection*{.getKeyNameExdr}
\subsubsection*{.getPaginateTarget}
\subsubsection*{.getPaginate}


\section{OfbizCatalogAltUrlTransform}
	\subsection{Functional role}
	\subsection{Application of the checklist}
\begin{itemize}
	\item[23.]{The javadoc is not used as shows at this link:\\
https://ci.apache.org/projects/ofbiz/site/javadocs/org/apache/ofbiz/product/category/ftl/OfbizCatalogAltUrlTransform.html}
	\item[25.]{line 50\\
There is no class comment. It must be added.
}
\end{itemize}
\subsubsection*{.getStringArg}
\subsubsection*{.checkArg}
\subsubsection*{.getWriter}
\begin{itemize}
 \item[various.]{line 143
		\begin{lstlisting} 
out.write(url.toString());
		\end{lstlisting}
		url is declared as String and as the .toString() javadoc for String type suggests 
\begin{displayquote}
		String java.lang.String.toString()

This object (which is already a string!) is itself returned.
\end{displayquote}}
\end{itemize}





\clearpage
\section{Effort Spent} %In this section you will include information about the number of hours each group member  has worked towards the fulfillment of this deadline.
	\begin{itemize}
		\item Alessandro Erba $\approx$ 19h
		\item Filippo Leveni 	$\approx$ 17h
		\item Luca Lodi $\approx$ 15h
	\end{itemize}

\end{document}